\chapter{\texorpdfstring{$\beta$}{Beta}-eliminação}\label{ch:b1:elimination}

Trate o 2-bromo-2-metilpropano com etóxido de sódio em etanol à
temperatura ambiente e o produto principal é um éter; aqueça a mesma
mistura e o produto principal é um alceno, o 2-metilpropeno, sem nenhum
grupo etóxi. O íon etóxido agiu não como \omterm{def:b1:electronic-effects:nucleophile}{nucleófilo}, atacando o carbono,
mas como base, retirando um próton do carbono vizinho enquanto o brometo
saía. Essa competição entre substituição e eliminação atravessa toda a
síntese orgânica. Este capítulo descreve os dois mecanismos de
eliminação, sua notável exigência geométrica, a regra que prevê qual
alceno se forma e como orientar uma reação em um sentido ou no outro.

\begin{recall}
SN1 e SN2 (\cref{ch:b1:nucleophilic-substitution}); \omterm{def:b1:stereochemistry-in-depth:representations}{projeções de Newman},
\omterm{def:b1:stereochemistry-in-depth:isomers}{conformações} do ciclo-hexano e descritores $Z$/$E$
(\cref{ch:b1:stereochemistry-in-depth}); \omterm{def:b1:electronic-effects:intermediates}{carbocátions} e bases
(\cref{ch:b1:electronic-effects}).
\end{recall}

\section{O mecanismo E2}

\begin{definition}[\texorpdfstring{$\beta$}{Beta}-eliminação, E2 e E1]\label{def:b1:elimination:elimination}
Uma \emph{$\beta$-eliminação}\index{beta-eliminacao@$\beta$-eliminação} retira um \omterm{def:b1:electronic-effects:nucleophile}{grupo de saída}
Y de um carbono (o carbono $\alpha$) e um próton de um carbono vizinho
(um carbono $\beta$), formando uma ligação dupla entre eles. No
\emph{mecanismo E2}\index{mecanismo E2}, uma base retira o próton enquanto
Y sai, em uma única \omterm{def:b1:elementary-steps:elementary-step}{etapa elementar}; no \emph{mecanismo
E1}\index{mecanismo E1}, Y sai primeiro, dando um \omterm{def:b1:electronic-effects:intermediates}{carbocátion}, que perde
em seguida um próton $\beta$.
\end{definition}

\begin{proposition}[Lei de velocidade da E2]\label{prop:b1:elimination:e2-law}
Para a E2, $v = k[\mathrm{RY}][\mathrm{B}]$, em que B é a base.
\end{proposition}

\begin{proof}
Uma única \omterm{def:b1:elementary-steps:elementary-step}{etapa elementar} bimolecular, envolvendo o \omterm{def:b1:nucleophilic-substitution:substitution}{substrato} e a base:
a lei da ação das massas das \omterm{def:b1:elementary-steps:elementary-step}{etapas elementares}
(\cref{ch:b1:elementary-steps}).
\end{proof}

\begin{definition}[Antiperiplanar e sinperiplanar]\label{def:b1:elimination:periplanar}
Duas ligações em átomos vizinhos são \emph{antiperiplanares}\index{antiperiplanar}
quando ficam em um mesmo plano, em lados opostos (\omterm{def:b1:stereochemistry-in-depth:conformations}{ângulo de torção} de
$180^\circ$), e \emph{sinperiplanares}\index{sinperiplanar} quando ficam em um
mesmo plano, do mesmo lado ($0^\circ$).
\end{definition}

\begin{omfigure}
\begin{tikzpicture}[font=\small, >={Stealth}]
  \pic (n) at (0,0) {omnewman={90}{270}};
  \node at (n-f1) {H}; \node at (n-f2) {\ce{CH3}}; \node at (n-f3) {H};
  \node at (n-b1) {Br}; \node at (n-b2) {\ce{CH3}}; \node at (n-b3) {H};
  \node (b) at (-2.6,2.0) {\ce{EtO-}};
  \draw[->, thick, omThm] (b) to[bend left=20] ($(n-f1)+(-0.15,0.1)$);
  \draw[->, thick, omThm] (0.12,0.8) to[bend left=40] (0.12,0.25);
  \draw[->, thick, omThm] (0.12,-0.6) to[bend right=35] ($(n-b1)+(0.25,0.15)$);
  \node[below] at (0,-2.0) {H e Br antiperiplanares};
  \draw[->, thick] (2.4,0) -- (3.6,0) node[midway, above] {E2};
  \node at (5.6,0) {\chemfig{C(-[:150]H_3C)(-[:210]H)=C(-[:30]H)(-[:-30]CH_3)}};
  \node at (5.6,-0.9) {\cip{E}-but-2-eno};
\end{tikzpicture}

{\small Reação E2 do 2-bromobutano, vista ao longo da ligação C3--C2 (C3
na frente, C2, que leva o Br, atrás) em uma de suas duas \omterm{def:b1:stereochemistry-in-depth:isomers}{conformações}
\omterm{def:b1:elimination:periplanar}{antiperiplanares}; a outra, com o outro hidrogênio de C3 em anti ao Br,
dá o \cip{Z}-but-2-eno. Três pares de elétrons se movem juntos: a base
retira o próton, o par C--H vira a ligação $\pi$, o par C--Br sai com o
bromo. Os grupos que ficam frente a frente na \omterm{def:b1:stereochemistry-in-depth:representations}{projeção de Newman} terminam
do mesmo lado da ligação dupla.}
\end{omfigure}

\begin{proposition}[A E2 é anti e estereoespecífica]\label{prop:b1:elimination:anti}
A E2 ocorre a partir da \omterm{def:b1:stereochemistry-in-depth:isomers}{conformação} em que as ligações C--H e C--Y são
\omterm{def:b1:elimination:periplanar}{antiperiplanares}. Ela é estereoespecífica: \omterm{def:b1:nucleophilic-substitution:substitution}{substratos} diastereoisoméricos
dão \omterm{def:b1:stereochemistry-in-depth:isomers}{estereoisômeros} diferentes do alceno, fixados pelo arranjo anti.
\end{proposition}

\begin{proof}
No \omterm{def:b1:elementary-steps:energy-profile}{estado de transição}, o \omterm{def:b1:lewis-resonance-vsepr:lewis}{par ligante} C--H precisa entrar na região
antiligante da ligação C--Y enquanto os dois carbonos se tornam
trigonais; a sobreposição é máxima quando as duas ligações são
paralelas, em um mesmo plano. O arranjo anti é alternado (de baixa
energia) e mantém a base longe do \omterm{def:b1:electronic-effects:nucleophile}{grupo de saída}; o arranjo sin é
eclipsado. Com H e Y em anti, os dois outros grupos de cada carbono
ficam fixados: os que estão do mesmo lado na \omterm{def:b1:stereochemistry-in-depth:representations}{projeção de Newman} terminam
em \textit{cis} na ligação dupla. Um \omterm{def:b1:stereochemistry-in-depth:enantiomers}{diastereoisômero} tem dois grupos
trocados em um carbono e dá o outro alceno.
\end{proof}

\begin{proposition}[A E2 em um ciclo-hexano]\label{prop:b1:elimination:cyclohexane-e2}
Em um anel de ciclo-hexano, a E2 exige que o \omterm{def:b1:electronic-effects:nucleophile}{grupo de saída} e o
hidrogênio $\beta$ sejam ambos \omterm{def:b1:stereochemistry-in-depth:conformations}{axiais} (\textit{trans}-diaxiais); um \omterm{def:b1:electronic-effects:nucleophile}{grupo
de saída} mantido \omterm{def:b1:stereochemistry-in-depth:conformations}{equatorial} não pode reagir até que o anel se inverta.
\end{proposition}

\begin{proof}
Em carbonos adjacentes de uma \omterm{def:b1:stereochemistry-in-depth:conformations}{cadeira}, duas ligações \omterm{def:b1:stereochemistry-in-depth:conformations}{axiais} são
\omterm{def:b1:elimination:periplanar}{antiperiplanares} (uma para cima, outra para baixo, paralelas ao eixo);
uma ligação \omterm{def:b1:stereochemistry-in-depth:conformations}{equatorial} só fica em anti com ligações C--C do anel, nunca
com uma ligação C--H.
\end{proof}

\begin{omfigure}
\begin{tikzpicture}[font=\small]
  \pic (a) at (0,0) {omchair};
  \draw[thick, omThm] (a-c1) -- (a-a1) node[above] {Cl};
  \draw[thick] (a-c1) -- (a-e1) node[right] {H};
  \draw[thick, omDef] (a-c2) -- (a-a2) node[below] {H};
  \draw[thick, omDef] (a-c6) -- (a-a6) node[below] {H};
  \node[align=left, anchor=west] at (3.0,0.2) {Cl axial (para cima) em C1;\\H axial (para baixo) nos dois vizinhos:\\dois H antiperiplanares};
\end{tikzpicture}

{\small Arranjo \textit{trans}-diaxial em uma \omterm{def:b1:stereochemistry-in-depth:conformations}{cadeira}: o cloro \omterm{def:b1:stereochemistry-in-depth:conformations}{axial} é
\omterm{def:b1:elimination:periplanar}{antiperiplanar} aos hidrogênios \omterm{def:b1:stereochemistry-in-depth:conformations}{axiais} dos dois carbonos vizinhos (em
azul).}
\end{omfigure}

\section{O mecanismo E1}

\begin{proposition}[Lei de velocidade da E1]\label{prop:b1:elimination:e1-law}
Para a E1, $v = k[\mathrm{RY}]$, sendo a velocidade a da ionização.
\end{proposition}

\begin{proof}
Como para a SN1 (\cref{prop:b1:nucleophilic-substitution:sn1-law}): o
\omterm{def:b1:electronic-effects:intermediates}{carbocátion} se forma em uma etapa lenta e é consumido depressa, aqui pela
perda de um próton para o solvente ou para uma \omterm{def:b1:predominant-reaction:strong-acid}{base fraca}; a
\omterm{def:b1:elementary-steps:qssa}{aproximação do estado estacionário} deixa $v = k_1[\mathrm{RY}]$.
\end{proof}

A E1 partilha seu \omterm{def:b1:electronic-effects:intermediates}{carbocátion} com a SN1: os \omterm{def:b1:nucleophilic-substitution:substitution}{substratos} terciários em
\omterm{def:b1:intermolecular-forces-solvents:solvent-classes}{solventes próticos} dão os dois produtos, e o alceno é favorecido pelo
aquecimento. Como o \omterm{def:b1:electronic-effects:intermediates}{carbocátion} é plano e o próton só é perdido depois,
a E1 não tem exigência geométrica e não é estereoespecífica; ela dá
principalmente o alceno mais estável, em geral o \cip{E}.
\begin{omfigure}
\setchemfig{atom sep=2.1em}
\schemestart
\chemfig{C(-[2]CH_3)(-[4]H_3C)(-[6]CH_3)-Br}
\arrow{->[lenta]}
\chemfig{C^{+}(-[2]CH_3)(-[:210]H_3C)(-[:-30]CH_3)}
\arrow{->[rápida][\ce{-H+}]}
\chemfig{C(-[2]CH_3)(-[:210]H_3C)=[:-30]CH_2}
\schemestop
\par\vspace{4pt}

{\small E1: ionização do 2-bromo-2-metilpropano ao \omterm{def:b1:electronic-effects:intermediates}{carbocátion} e depois
perda de um próton de um grupo metila para o solvente, dando o
2-metilpropeno.}
\end{omfigure}

\section{A regiosseletividade: a regra de Zaitsev}

\begin{definition}[Reações regiosseletivas e estereosseletivas]\label{def:b1:elimination:selectivity}
Uma reação que pode dar vários \omterm{def:b1:stereochemistry-in-depth:isomers}{isômeros constitucionais}, mas forma
principalmente um deles, é uma \emph{reação regiosseletiva}\index{reacao regiosseletiva@reação regiosseletiva};
uma reação que pode dar vários \omterm{def:b1:stereochemistry-in-depth:isomers}{estereoisômeros}, mas forma principalmente
um deles, é uma \emph{reação estereosseletiva}\index{reacao estereosseletiva@reação estereosseletiva}.
\end{definition}

\begin{proposition}[Regra de Zaitsev]\label{prop:b1:elimination:zaitsev}
Quando vários hidrogênios $\beta$ podem ser retirados, o alceno
majoritário é em geral o mais substituído (o mais estável), e o isômero
\cip{E} de preferência ao \cip{Z} (\emph{regra de Zaitsev}\index{regra de Zaitsev}).
Bases volumosas e restrições geométricas, como a exigência
trans-diaxial, podem contrariá-la.
\end{proposition}

\begin{proof}
Os grupos alquila nos carbonos de uma ligação dupla a estabilizam (como
estabilizam os \omterm{def:b1:electronic-effects:intermediates}{carbocátions}), e os alcenos \cip{E} evitam o
congestionamento dos grupos em \textit{cis}. Os \omterm{def:b1:elementary-steps:energy-profile}{estados de transição} da
E2 e da E1 têm caráter parcial de ligação dupla, de modo que o alceno
mais estável se forma mais depressa. Uma base volumosa alcança mais
facilmente os hidrogênios menos impedidos; e um hidrogênio que não pode
ficar \omterm{def:b1:elimination:periplanar}{antiperiplanar} ao \omterm{def:b1:electronic-effects:nucleophile}{grupo de saída} não é retirado pela E2, qualquer
que seja a estabilidade do alceno que daria.
\end{proof}

\begin{method}[Prever o alceno]\label{met:b1:elimination:predict-alkene}
\begin{enumerate}
  \item Liste os carbonos $\beta$ que levam hidrogênios.
  \item Para a E2, conserve apenas os hidrogênios que podem ficar
        \omterm{def:b1:elimination:periplanar}{antiperiplanares} a Y (em um anel: trans-diaxiais, em uma \omterm{def:b1:stereochemistry-in-depth:conformations}{cadeira}
        que possa de fato se formar).
  \item Entre eles, favoreça o que dá o alceno mais substituído
        (Zaitsev), a menos que a base seja volumosa.
  \item Para cada um, desenhe a \omterm{def:b1:stereochemistry-in-depth:representations}{projeção de Newman} na \omterm{def:b1:stereochemistry-in-depth:conformations}{conformação anti}
        para ler a geometria \cip{E}/\cip{Z}.
\end{enumerate}
\end{method}

\section{Substituição ou eliminação?}

\begin{proposition}[Substituição contra eliminação]\label{prop:b1:elimination:sn-vs-e}
A eliminação é favorecida pelas \omterm{def:b1:predominant-reaction:strong-acid}{bases fortes} e volumosas, pela
temperatura alta e pelos \omterm{def:b1:nucleophilic-substitution:substitution}{substratos} terciários ou congestionados; a
substituição, pelos bons \omterm{def:b1:electronic-effects:nucleophile}{nucleófilos} que são \omterm{def:b1:predominant-reaction:strong-acid}{bases fracas}, pela
temperatura baixa e pelos \omterm{def:b1:nucleophilic-substitution:substitution}{substratos} primários.
\end{proposition}

\begin{proof}
A força da base favorece o ataque ao hidrogênio; o volume dificulta o
ataque a um carbono congestionado, mas não a um hidrogênio exposto. Um
\omterm{def:b1:electronic-effects:carbon-class}{carbono terciário} é inacessível à SN2, e seu \omterm{def:b1:electronic-effects:intermediates}{carbocátion} perde prótons
com facilidade. A eliminação transforma uma molécula em duas ou três
(alceno, \omterm{def:b1:electronic-effects:nucleophile}{grupo de saída}, base protonada): ela é favorecida quando a
temperatura sobe, um resultado justificado no volume do 2.º~ano.
\end{proof}

\begin{omfigure}
\begin{tabular}{llll}
\hline
\omterm{def:b1:nucleophilic-substitution:substitution}{substrato} & \omterm{def:b1:predominant-reaction:strong-acid}{base forte}, pequena & \omterm{def:b1:predominant-reaction:strong-acid}{base forte}, volumosa & \omterm{def:b1:predominant-reaction:strong-acid}{base fraca}, bom \omterm{def:b1:electronic-effects:nucleophile}{nucleófilo} \\
\hline
primário & SN2 (um pouco de E2) & E2 & SN2 \\
secundário & E2 (e SN2) & E2 & SN2 (aprótico), SN1/E1 (prótico) \\
terciário & E2 & E2 & SN1/E1 (prótico) \\
\hline
\end{tabular}

\smallskip
{\small Uma primeira orientação: a reação principal em cada caso. O
aquecimento favorece a eliminação, coluna por coluna.}
\end{omfigure}

\section{Exercícios}

\begin{exercise}[$\star$]\label{exo:b1:elimination:1}
Dê os alcenos que podem se formar por eliminação de HBr a partir do
2-bromobutano, do 2-bromo-2-metilbutano e do bromociclo-hexano, e aquele
que a \omterm{prop:b1:elimination:zaitsev}{regra de Zaitsev} prevê como majoritário.
\end{exercise}

\begin{exercise}[$\star$]\label{exo:b1:elimination:2}
Escreva as \omterm{def:b1:rate-laws:order}{leis de velocidade} da E1 e da E2. Como você as distinguiria
experimentalmente para o 2-bromo-2-metilpropano em etanol contendo
etóxido?
\end{exercise}

\begin{exercise}[$\star$]\label{exo:b1:elimination:3}
Desenhe as \omterm{def:b1:stereochemistry-in-depth:representations}{projeções de Newman} do 2-bromobutano ao longo de C2--C3 em
suas três \omterm{def:b1:stereochemistry-in-depth:conformations}{conformações alternadas} e diga quais podem sofrer E2.
\end{exercise}

\begin{exercise}[$\star$]\label{exo:b1:elimination:4}
Preveja a reação principal (SN1, SN2, E1, E2): bromoetano com hidróxido
de sódio em água a \qty{25}{\celsius}; 2-bromo-2-metilpropano com
terc-butóxido de potássio; 2-bromopropano com iodeto de sódio em
propanona; 2-bromo-2-metilpropano em etanol quente.
\end{exercise}

\begin{exercise}[$\star\star$]\label{exo:b1:elimination:5}
Mostre, com \omterm{def:b1:stereochemistry-in-depth:representations}{projeções de Newman}, que a eliminação E2 de HBr a partir dos
dois \omterm{def:b1:stereochemistry-in-depth:enantiomers}{diastereoisômeros} do 2-bromo-3-metilpentano que podem perder o
hidrogênio de C3 dá os dois \omterm{def:b1:stereochemistry-in-depth:isomers}{estereoisômeros} do 3-metilpent-2-eno.
\end{exercise}

\begin{exercise}[$\star\star$]\label{exo:b1:elimination:6}
Escreva o mecanismo da desidratação do 2-metilbutan-2-ol em ácido
sulfúrico concentrado quente (E1) e dê o alceno majoritário.
\end{exercise}

\begin{exercise}[$\star\star$]\label{exo:b1:elimination:7}
Explique por que o \textit{trans}-1-bromo-4-terc-butilciclo-hexano reage
muito lentamente por E2, enquanto seu isômero \textit{cis} reage depressa.
(O grupo terc-butila permanece \omterm{def:b1:stereochemistry-in-depth:conformations}{equatorial}.)
\end{exercise}

\begin{exercise}[$\star\star$]\label{exo:b1:elimination:8}
O 2-bromo-2-metilbutano com etóxido de sódio dá principalmente
2-metilbut-2-eno; com terc-butóxido de potássio, principalmente
2-metilbut-1-eno. Explique.
\end{exercise}

\begin{exercise}[$\star\star$]\label{exo:b1:elimination:9}
Por que nunca se observa eliminação com o bromometano? E com o
1-bromo-2,2-dimetilpropano por E2?
\end{exercise}

\begin{exercise}[$\star\star\star$]\label{exo:b1:elimination:10}
Em etanol a \qty{25}{\celsius}, o 2-bromo-2-metilpropano dá uma mistura
do éter (SN1) e do alceno (E1) cujas proporções não dependem da
concentração do \omterm{def:b1:nucleophilic-substitution:substitution}{substrato}. Mostre que os dois produtos vêm do mesmo
\omterm{def:b1:electronic-effects:intermediates}{carbocátion} e que sua razão é a razão entre duas \omterm{def:b1:rate-laws:order}{constantes de
velocidade}.
\end{exercise}

\begin{exercise}[$\star\star\star$]\label{exo:b1:elimination:11}
Um \omterm{def:b1:stereochemistry-in-depth:enantiomers}{diastereoisômero} do 1,2-dibromo-1,2-difeniletano (o meso) dá, por E2
com etóxido, um único \omterm{def:b1:stereochemistry-in-depth:isomers}{estereoisômero} do 1-bromo-1,2-difenileteno.
Encontre qual, com uma \omterm{def:b1:stereochemistry-in-depth:representations}{projeção de Newman}.
\end{exercise}

\begin{exercise}[$\star\star\star$]\label{exo:b1:elimination:12}
Usando as energias do \cref{ch:b1:stereochemistry-in-depth}, explique por
que uma reação E2 em um ciclo-hexano pode ser desacelerada por um fator
de centenas quando o \omterm{def:b1:electronic-effects:nucleophile}{grupo de saída} precisa ficar \omterm{def:b1:stereochemistry-in-depth:conformations}{axial} em uma \omterm{def:b1:stereochemistry-in-depth:conformations}{cadeira}
congestionada. Expresse a velocidade como o produto de uma fração de
confôrmero por uma \omterm{def:b1:rate-laws:order}{constante de velocidade}.
\end{exercise}

\section{Problema: Os cloretos de mentila e de neomentila}

\begin{problem}[{Problema de fim de semana --- as cadeiras de dois cloretos
diastereoisoméricos, os hidrogênios antiperiplanares disponíveis em cada
um, os alcenos formados e a razão entre suas velocidades de E2}]
\label{pb:b1:elimination:1}
O cloreto de mentila e o cloreto de neomentila são os cloretos do mentol
e do neomentol (\cref{ch:b1:stereochemistry-in-depth}): no anel, C1 leva
o Cl, C2 o grupo isopropila, C5 o grupo metila. No cloreto de mentila, os
três grupos podem ficar todos \omterm{def:b1:stereochemistry-in-depth:conformations}{equatoriais}; no cloreto de neomentila, o Cl
é \omterm{def:b1:stereochemistry-in-depth:conformations}{axial} quando os dois grupos alquila são \omterm{def:b1:stereochemistry-in-depth:conformations}{equatoriais}. Ambos são
tratados com etóxido de sódio em etanol, uma reação E2. Dados do
problema: \qty{95}{\percent} do cloreto de neomentila está na \omterm{def:b1:stereochemistry-in-depth:conformations}{cadeira} com
Cl \omterm{def:b1:stereochemistry-in-depth:conformations}{axial}; para o cloreto de mentila, a \omterm{def:b1:stereochemistry-in-depth:conformations}{cadeira} com Cl \omterm{def:b1:stereochemistry-in-depth:conformations}{axial} (os três
grupos \omterm{def:b1:stereochemistry-in-depth:conformations}{axiais}) é uma fração de \num{5.0e-3}; suponha que cada hidrogênio
\omterm{def:b1:elimination:periplanar}{antiperiplanar} de uma \omterm{def:b1:stereochemistry-in-depth:conformations}{cadeira} reativa reage com a mesma \omterm{def:b1:rate-laws:order}{constante de
velocidade}. O cloreto de neomentila dá \qty{78}{\percent} do alceno
trissubstituído.

\medskip
\noindent\textbf{Parte I --- As \omterm{def:b1:stereochemistry-in-depth:conformations}{cadeiras}.}
\begin{enumerate}
  \item Desenhe a \omterm{def:b1:stereochemistry-in-depth:conformations}{cadeira} preferida do cloreto de mentila.
  \item Desenhe a \omterm{def:b1:stereochemistry-in-depth:conformations}{cadeira} preferida do cloreto de neomentila.
  \item Os dois cloretos são \omterm{def:b1:stereochemistry-in-depth:enantiomers}{enantiômeros} ou \omterm{def:b1:stereochemistry-in-depth:enantiomers}{diastereoisômeros}?
  \item Enuncie a exigência geométrica da E2 em um anel.
  \item Em que \omterm{def:b1:stereochemistry-in-depth:conformations}{cadeira} o cloreto de mentila pode reagir? Desenhe-a.
  \item Por que essa \omterm{def:b1:stereochemistry-in-depth:conformations}{cadeira} é rara? Use o custo energético dos grupos
        \omterm{def:b1:stereochemistry-in-depth:conformations}{axiais}.
\end{enumerate}

\medskip
\noindent\textbf{Parte II --- Os hidrogênios anti.}
\begin{enumerate}[resume]
  \item Liste os carbonos $\beta$ do cloreto (C2 e C6) e seus
        hidrogênios.
  \item Na \omterm{def:b1:stereochemistry-in-depth:conformations}{cadeira} reativa do cloreto de neomentila, quais hidrogênios
        são \omterm{def:b1:elimination:periplanar}{antiperiplanares} ao Cl?
  \item Na \omterm{def:b1:stereochemistry-in-depth:conformations}{cadeira} reativa do cloreto de mentila, o hidrogênio de C2 é
        \omterm{def:b1:stereochemistry-in-depth:conformations}{axial}? Um hidrogênio de C6 é \omterm{def:b1:stereochemistry-in-depth:conformations}{axial}?
  \item Quantos hidrogênios \omterm{def:b1:elimination:periplanar}{antiperiplanares} cada cloreto oferece?
  \item Desenhe uma \omterm{def:b1:stereochemistry-in-depth:representations}{projeção de Newman} ao longo de C1--C6 para a \omterm{def:b1:stereochemistry-in-depth:conformations}{cadeira}
        reativa do cloreto de mentila.
  \item Por que uma eliminação \omterm{def:b1:elimination:periplanar}{sinperiplanar} não poderia salvar a
        reação?
\end{enumerate}

\medskip
\noindent\textbf{Parte III --- Os produtos.}
\begin{enumerate}[resume]
  \item Dê o nome do alceno formado pela retirada do hidrogênio de C2 e o
        do formado pela retirada de um hidrogênio de C6.
  \item Qual deles a \omterm{prop:b1:elimination:zaitsev}{regra de Zaitsev} favorece, e por quê?
  \item O que dá o cloreto de neomentila? A razão 78/22 é coerente com a
        \omterm{prop:b1:elimination:zaitsev}{regra de Zaitsev}?
  \item O que dá o cloreto de mentila? Isso é coerente com a \omterm{prop:b1:elimination:zaitsev}{regra de
        Zaitsev}?
  \item Que reação é estereoespecífica, regiosseletiva, ou ambas?
  \item Por que a substituição (SN2) é minoritária com o etóxido aqui?
\end{enumerate}

\medskip
\noindent\textbf{Parte IV --- As velocidades.}
\begin{enumerate}[resume]
  \item Escreva a velocidade de cada reação em função da fração de
        \omterm{def:b1:stereochemistry-in-depth:conformations}{cadeira} reativa e do número de hidrogênios \omterm{def:b1:elimination:periplanar}{antiperiplanares}.
  \item Calcule a velocidade relativa do cloreto de neomentila.
  \item Calcule a velocidade relativa do cloreto de mentila.
  \item Calcule a razão entre as velocidades $k(\text{neomentila})/k(\text{mentila})$.
\end{enumerate}
\end{problem}
